\chapter{Introduction}

\section{Invitation}

The study of physics since its modern formulation in the times of Newton
and Galileo has found ever greater explanatory power in finding unifying
principles that relate seemingly disparate phenomena. While this insight
is not novel, the discovery of \emph{dualities} between physical
theories is a more recent take on the paradigm of unification, which has
spurred several lines of inquiry in theoretical physics. 

One such duality that arose from the study of amplitudes is \textbf{color-kinematics duality}, also known as \textbf{BCJ
duality} \cite{Bern:2008qj}\footnote{The two names are used interchangeably throughout this dissertation.}. BCJ
duality shows up as a relationship between amplitudes in Yang-Mills theories and gravitational theories --- in a
Feynman diagram expansion of Yang-Mills amplitudes, performing a term-by-term replacement of factors that
carry all the color information by a kinematic factor results in a gravitational amplitude with the same external
kinematics.

For instance, at tree level, a four-point scattering amplitude in pure Yang-Mills theory can be written as
\cite{Bern:2008qj}
\begin{equation} \label{intro:eq:YM4point} 
    \mathcal{A}^{\text{tree}}_{4} = g^{2} \left( \frac{n_{s}c_{s}}{s} + \frac{n_{t}c_{t}}{t} +
    \frac{n_{u}c_{u}}{u}\right),
\end{equation}
where the three terms represent contributions from the $s$-, $t$-, and $u$-channel diagrams
respectively.\footnote{In the usual Feynman diagram expansion, there is another diagram from the four-point vertex
    of a non-Abelian Yang-Mills theory. It turns out that the fourth diagram can be written as the sum of three
terms, which can be absorbed appropriately into the three terms above.} The
denominators $s,t,u$ are the usual Mandelstam variables, $s = (p_{1} + p_{2})^{2}, t = (p_{1} + p_{4})^{2}, u =
(p_{1} + p_{3})^{2}$, where $p_{i}$ are the external gluon momenta. The terms in the numerator are split into
\textit{color factors} $c_{i}$, which depend only on the Yang-Mills symmetry group, and \textit{kinematic factors}
$n_{i}$, which depend only on external momenta and polarizations.

A detailed account of the expressions $c_{i}$ and $n_{i}$ could be given here, but since the primary focus of this
dissertation will be on the ramifications of BCJ duality on classical solutions, we will state only that they
satisfy the identities,
\begin{equation} \label{} 
    c_{s} + c_{t} + c_{u} = 0, 
\end{equation}
which follows from the Jacobi identity of the Yang-Mills group, and 
\begin{equation} \label{} 
    n_{s} + n_{t} + n_{u} = 0,
\end{equation}
which follows from putting external momenta on-shell, and is sometimes termed the \textit{kinematic Jacobi
identity}\footnote{It is the existence of such kinematic factors, which satisfy all corresponding color factor
identities, that leads to the proposal of a duality between color and kinematics.}. Hence, performing the replacement $c_{i} \to n_{i}$, along with replacing the gauge theory coupling,
leads to the tree-level four point amplitude of a different field theory,
\begin{equation} \label{} 
    \mathcal{M}^{\text{tree}}_{4} = \kappa^{2} \left( \frac{n_{s}^{2}}{s} + \frac{n_{t}^{2}}{t} +
    \frac{n_{u}^{2}}{u}\right).
\end{equation}
This is in fact the tree level amplitude for four-point graviton scattering. Since the amplitude has two copies of
the kinematic factor in the numerator, it is called the \textit{double copy} of the Yang-Mills amplitude
\eqref{intro:eq:YM4point}. It is also customary to refer to gravitational theory itself as the \textit{double copy}
of the Yang-Mills theory, and to refer to the Yang-Mills theory as the \textit{single copy}. The duality between color and kinematics can also be applied in the other direction, $n_{i} \to \tilde{c}_{i}$,
\begin{equation} \label{} 
    \mathcal{A}_{4}^{\text{tree}} \to y^{2}\left( \frac{c_{s}\tilde{c}_{s}}{s} + \frac{c_{t}\tilde{c}_{t}}{t} +
    \frac{c_{u}\tilde{c}_{u}}{u} \right),
\end{equation}
which results in the four-point amplitude for a massless \textit{biadjoint scalar field}. A biadjoint scalar
transforms in the adjoint representation of two (possibly different) gauge groups, with coupling constant $y$.
This biadjoint scalar theory is usually dubbed the \textit{zeroth copy}.

Such double copy relations between gauge and gravity amplitudes have been proven at tree level \cite{Bern:2010yg},
and are conjectured to hold to all loop orders, with an ever-expanding set of examples (see, for instance, chapter
6 of the review \cite{Bern:2019prr}). The double copy, therefore, has shown us how to construct amplitudes for a wide
range of theories (scalar, gauge, gravity) from a common set of building blocks in a way that does not obviously
follow from the Lagrangians themselves. While our interest in the double copy paradigm was motivated by probing the
extent of the double copy proposal and hopefully learn something fundamental about the nature of gravitational
theories, the double copy has also given rise to applications to other fields of physics \cite{Adamo:2022dcm}.

\section{From Quantum Amplitudes to Classical Physics}

%\begin{itemize}
    %\item  Duality extends to
%classical physics --- the double copy. 
    %\item  Perturbative double copy ==
%color-kinematics of tree level 
%\item  \st{Exact solution approaches:} 
    %\begin{itemize}
        %\item \st{Kerr-Schild} 
        %\item  \st{Weyl} 
        %\item \st{Newman-Penrose} 
        %\item  \st{Twistors --- in cases where linearized solutions are enough to construct exact solution}
    %\end{itemize}
%\item  Why study the classical limit? 
%\item GW inspiral waveforms can be recast
%as classical limits of quantum amplitudes 
%\item Other gravitational observables from (limits of?) on-shell quantum amplitudes cf. Kosower, Maybee, O'Connell
    %\href{https://www.arxiv.org/abs/1811.10950}{1811.10950} 
%\item New tool for
%systematic study of exact solutions in gravity 
%\item Curiosity
%\end{itemize}

The discovery of a double copy structure in the amplitudes computed in gauge and gravity theories naturally leads
to a similar structure in their respective classical limits. Since the Feynman diagram series with an increasing
number of loops is effectively an expansion in increasing powers of $\hbar$, taking the classical limit of an
amplitude generally corresponds to keeping only the tree-level terms and discarding
all diagrams with loops. See the recent review \cite{Kosower:2022yvp} for a comprehensive description of how the
classical double copy emerges from the double copy of amplitudes.

The study of the classical double copy via amplitudes expanded beyond
the amplitudes perspective and took on a life of its own after a seminal paper \cite{Monteiro:2014cda}
that described how a certain class of spacetimes (Kerr-Schild spacetimes) could be constructed
\emph{exactly} from solutions in the corresponding electromagnetic and
scalar theories. Kerr-Schild solutions are distinguished by the ability of the
metric to be written as 
\begin{equation} \label{} 
g_{\mu \nu} = \eta_{\mu \nu} + \phi k_{\mu} k_{\nu},
\end{equation}
where $k_{\mu}$ is a null vector field tangent everywhere to affinely parametrized geodesics, $\phi$ is a scalar
function, and $\eta_{\mu \nu}$ is the Minkowski metric. The authors of \cite{Monteiro:2014cda} showed that the corresponding single and zeroth copies are,
\begin{equation} \label{Intro:Kerr-SchildSingleZerothCopies} 
    A_{\mu} = \phi k_{\mu}, \qquad \Phi = \phi,
\end{equation}
which satisfy equations of motion (for a $U(1)$ gauge field and a massless scalar field) when considered to be
fields on Minkowski spacetime with metric $\eta_{\mu \nu}$. Since this proposal of the \textit{Kerr-Schild double
copy}, there have been proposals for other formulations that extend the classical double copy in various ways. The
most notable of these for our purposes is the \textit{Weyl double copy} \cite{Luna:2018dpt, Godazgar:2020zbv}.

This dissertation is a collection of our attempts to extend the double copy in various directions, and to probe the
limits of its applicability. The following section gives an outline of the subsequent chapters, which each tackle a
specific way of extending the double copy

\section{Extensions}

\subsection{Higher Dimensions}
%\begin{itemize}
%\tightlist
%\item
  %Why study higher dimensions?

  %\begin{itemize}
  %\tightlist
  %\item
    %See {[}{[}Meeting with Cindy Jan 15\#Why higher dimensions?{]}{]}
  %\end{itemize}
%\item
    %\st{Kerr-Schild formalism works in arbitrary dimension, but (spinor) Weyl
  %needs dimension-specific spinor formalisms.}
  
%\item
  %\st{Showed that for Kerr-NUT-(A)dS, the previously discovered aligned
  %fields did in fact double copy to the spacetime when considered on an
%appropriate background.}
%\item
    %\st{Point to chapter 2.}
%\end{itemize}
Our first interest was in extending the classical double copy to higher dimensions. The Kerr-Schild formalism as
described above (eq \eqref{Intro:Kerr-SchildSingleZerothCopies}) generalizes easily to higher dimensions, which is
an advantage of the formalism. What we give up in
return is gauge invariance;
%demonstrating that the metric is of Kerr-Schild type is usually done via
%transforming to a specific coordinate system in which $\eta_{\mu \nu}$ takes the usual (Cartesian) form,
%\begin{equation} \label{} 
    %[\eta_{\mu \nu}] = 
    %\begin{pmatrix}
        %-1 & 0 & 0 & 0 \\
        %0 & 1 & 0 & 0 \\
        %0 & 0 & 1 & 0 \\
        %0 & 0 & 0 & 1
    %\end{pmatrix}.
%\end{equation}
%Arguably, this is not mandatory. It can in principle be shown in any coordinate system that the $\eta_{\mu \nu}$
%term represents a flat metric. 
%At the level of the single copy however,
the prescription for $A_{\mu}$ in \ref{Intro:Kerr-SchildSingleZerothCopies} makes an explicit gauge choice. After a
generic gauge transformation,
\begin{equation} \label{} 
    A_{\mu} \to A_{\mu} + \partial_{\mu} \alpha,
\end{equation}
$A_{\mu}$ can no longer be expressed as being proportional to a null geodesic vector field. 

The Weyl double copy proposed in \cite{Luna:2018dpt} remedies the gauge dependence of the Kerr-Schild
approach by working with the electromagnetic field strength $F_{\mu \nu}$ instead of the gauge field
$A_{\mu}$, and the Weyl tensor $W_{\mu \nu \rho \sigma}$ instead of the metric $g_{\mu \nu}$. It turns out that the
Weyl double copy is most naturally stated in terms of the spinor representation of these tensors, and in $D = 4$
the Weyl double copy is
\begin{equation} \label{} 
    \Psi_{ABCD} = \frac{1}{S} f_{(AB} f_{CD)},
\end{equation}
where $\Psi_{ABCD}$ is the Weyl spinor, $f_{AB}$ is the electromagnetic field strength, $S$ is a scalar related to
the zeroth copy, and $A, B, C, D = 1, 2$ are spinor indices.

While the gauge independence of this formalism is desirable, its reliance on a spinorial representation makes it
dimension-specific. In chapter \ref{chap:5D}, we explain a proposal for the five-dimensional version of the spinor Weyl
double copy. We demonstrate the double copy structure of the Kerr-NUT-(A)dS class of spacetimes --- in arbitrary
dimension using the Kerr-Schild double copy, and in five dimensions using our proposal for the spinor Weyl double
copy. We find that our results are consistent with previously known examples in the Kerr-Schild and Weyl spinor
formalisms.

\subsection{\texorpdfstring{Horizons }{Horizons }}

%\begin{itemize}
%\tightlist
%\item
  %If we had discovered double copy before classical GR, could we have
  %discovered black hole spacetimes (in particular their causal
  %horizons?)
%\item
  %Classical double copy formalisms (exact formulations) are local maps
  %=\textgreater{} First study local horizons i.e.~apparent horizons.
%\item
  %Kerr-Schild result in a remarkably simple map, but necessarily
  %involves the zeroth copy.
%\item
  %Also, the single copy appears in a gauge-dependent way. Not
  %necessarily bad since the double copy prescription also needs
  %gauge-fixing, but hard to interpret the horizon map.
%\item
  %Point to chapter 3.
%\end{itemize}

A rich area of inquiry within the double copy paradigm is building out the double copy dictionary beyond simply
calculating the gravitational, gauge, and scalar fields themselves. However, a question remained: how do black hole
horizons emerge from the double copy structure of black hole spacetimes? If the double copy proposal is to be taken seriously, in that gravitational physics is entirely constructible from
the same building blocks as the physics of a gauge field (plus an associated biadjoint scalar), there must be a way to
identify the presence of a horizon from the single and zeroth copy theories.

For instance, the single copy of the Schwarzschild black hole is the Coulomb potential of a static point charge,
and the zeroth copy (scalar field) is proportional to $1 / r$. From these fields, it is not immediately obvious how
the double copy's horizon at $r = GM$ could be diagnosed. In particular, the gauge potential and the scalar field
are smooth everywhere except $r = 0$.

We explore this question in chapter \ref{chap:Horizons}, where we give a proposal for finding trapped surfaces in
the gravitational solution using quantities that can be computed using the gauge field, the associated scalar, and
the supposition of a Kerr-Schild solution that is their double copy. The relevant quantities involve both the gauge
field and the scalar field together, but they are not invariant under gauge transformations of the gauge potential.

\subsection{\texorpdfstring{The Penrose limit
}{The Penrose limit }}

So far, we relied on double copy proposals for classical gravity that were both
\textit{exact}, and \textit{local} in position space. These methods are known to work for special classes of
spacetimes, one of which (Kerr-Schild metrics) we have already mentioned. The Weyl double copy in its exact form
requires the spacetime to be \textit{algebraically special}, which is a condition on the \textit{principal null
directions} of the Weyl tensor. More precise statements about algebraic speciality are made in \ref{ssec:5DTypeD}.

The pursuit of ever more general examples of the double copy remains a fruitful endeavor. However, recent work
\cite{Luna:2022dxo} showed that the two formalisms we have been working with (Kerr-Schild and Weyl) are expected to
fail for algebraically general spacetimes. More than a failure of exactness, the authors of \cite{Luna:2022dxo} set
out to answer why the classical double copy was local in the fields anyway --- the amplitudes double copy seems to
imply that the double copy should be local in momentum space, not in position space. It turns out that the
existence of a local double copy relies strongly on algebraic specialitay. 

It seems then that there is more to be gained from a perturbative approach. The perturbative classical double copy
arising from the amplitudes paradigm is naturally expressed as a perturbation about Minkowski spacetime. Now, with the
knowledge of exact classical double copies, we had the opportunity to deal with expansions about other background
spacetimes. One such approach was taken by \cite{Godazgar:2021iae}, where the Weyl spinor was expanded around null
infinity in an asymptotically flat spacetime. The Weyl spinor is algebraically special to leading orders in $1 /
r$, which allows for application of the Weyl double copy near null infinity.

Our motivation was to pursue a related line of inquiry: what happens if, instead of expanding about a null
boundary, we expand the Weyl spinor around a null
geodesic in the bulk? The limiting spacetime (the \textit{Penrose limit}) at any null geodesic in any spacetime was
shown to be a plane wave by Penrose \cite{penrose1976any}, which is an algebraically special spacetime.

Chapters \ref{chap:Penrose} and \ref{chap:PenroseExpansion} lay out our investigations in expanding about null
geodesics. Chapter \ref{chap:Penrose} is specifically about whether the process of taking the Penrose limit
commutes with the double copy map, which we expected on physical grounds but nevertheless had to explicitly check.
Chapter \ref{chap:PenroseExpansion} then examines how the double copy structure of a spacetime decomposes into
various terms in an expansion about the null geodesic. This allowed us to examine at what order in the Penrose
expansion the existence of a local double copy breaks down for generic spacetimes.

\nobibintoc
\printbibliography[heading=subbibliography]
